%%
%% The "author" command and its associated commands are used to define
%% the authors and their affiliations.
%% Of note is the shared affiliation of the first two authors, and the
%% "authornote" and "authornotemark" commands
%% used to denote shared contribution to the research.

\settopmatter{authorsperrow=2}
\author{Charlotte Kiesel}
% \authornote{Both authors contributed equally to this research.}
\email{yoder6@illinois.edu}
\affiliation{%
  \institution{Univ. of Illinois in Urbana-Champaign}
  % \city{Champaign}
  % \state{Illinois}
  \country{}
}
\orcid{0009-0006-7316-9351}

\author{Dipayan Mukherjee}
\email{dipayan2@illinois.edu}
\affiliation{%
  \institution{Univ. of Illinois in Urbana-Champaign}
  % \city{Champaign}
  % \state{Illinois}
  \country{}
}
\orcid{0000-0002-6017-1739}

\author{Mark Hasegawa-Johnson}
\email{jhasegaw@illinois.edu}
\affiliation{%
  \institution{Univ. of Illinois in Urbana-Champaign}
  % \city{Champaign}
  % \state{Illinois}
  \country{}
}
\orcid{0000-0002-5631-2893}

\author{Karrie Karahalios}
\email{kkarahal@illinois.edu}
\affiliation{%
  \institution{Univ. of Illinois in Urbana-Champaign}
  % \city{Champaign}
  % \state{Illinois}
  \country{}
}
\orcid{0000-0001-8788-3405}

%%
%% By default, the full list of authors will be used in the page
%% headers. Often, this list is too long, and will overlap
%% other information printed in the page headers. This command allows
%% the author to define a more concise list
%% of authors' names for this purpose.
\renewcommand{\shortauthors}{Kiesel et al.}
